
The goal is to test whether calibration quality, hallucination resistance, and adversarial robustness share a common latent factor — \textit{epistemic reliability} — independent of raw capability. The analysis treats models as subjects and benchmark scores as items, following the psychometric framing of Burnell et al. (2023) and Polo et al. (2024).

\subsubsection{3.1 The YouRA Framework}

The YouRA (Your Research Assistant) framework provides a multi-hypothesis evaluation pipeline that decomposes the main hypothesis into four sub-hypotheses:

\begin{itemize}
\item \textbf{H-E1 (Existence):} Cross-property partial Spearman correlations exceed |ρ| ≥ 0.40 with BCa CIs excluding zero, and factor analysis extracts a stable latent dimension (Tucker's congruence ≥ 0.85). Gate type: MUST_WORK.
\item \textbf{H-M1 (Calibration–Hallucination Mechanism):} The calibration–hallucination correlation survives MMLU capability control (survival fraction ≥ 0.50). Gate type: MUST_WORK.
\item \textbf{H-M2 (Predictive Power):} The epistemic composite predicts top-quartile adversarial failure beyond MMLU-only (ΔAUC ≥ 0.10, CI lower bound > 0). Gate type: SHOULD_WORK.
\item \textbf{H-M3 (Embedding Perturbation):} Calibration predicts decision-surface smoothness under Gaussian perturbation. Gate type: SHOULD_WORK. Pre-registered; not executed in this study.
\end{itemize}

\subsubsection{3.2 Model Population}

The target population is N=30 instruction-tuned open-weight LLMs, 7B–70B parameters, 8 families: LLaMA-2, Mistral, Falcon, Pythia, Qwen, Yi, OLMo, Gemma. Selected to maximize diversity of training regime, scale, and architecture. HuggingFace-accessible as of 2024-01. The current implementation uses a synthetic score matrix; this population is the target for real-data execution (FW1).

The score matrix used in this study includes models from the following families: llama2 (3 models), llama3 (2), mistral/mixtral (4), falcon (2), vicuna (2), zephyr (2), nous (2), wizardlm (2), qwen (3), yi (2), deepseek (2), internlm (2), phi (1), solar (1). The parameter range is 6.9B–70B.

\subsubsection{3.3 Benchmark Metrics}

\begin{table}[htbp]
\centering
\begin{tabular}{lll}
\toprule
Metric & Benchmark & Role \\
\midrule
ECE & MMLU logits (10-bin) & Calibration quality \\
Brier score & MMLU logits & Proper scoring calibration \\
TruthfulQA\% & TruthfulQA (817 items) & Hallucination resistance \\
AdvGLUE drop & AdvGLUE (5 tasks) & Adversarial robustness \\
ANLI drop & ANLI R1–R3 & Adversarial robustness \\
MMLU accuracy & MMLU (57 subjects) & Capability control variable only \\
\bottomrule
\end{tabular}
\end{table}

\subsubsection{3.4 Statistical Analysis}

\textbf{Partial Spearman correlation (H-E1, H-M1):} For each metric pair (X, Y), partial Spearman correlation is computed controlling for MMLU accuracy, with BCa bootstrap 95\% CIs (B=10,000, seed=42). Gate threshold: |ρ| ≥ 0.40, CI excludes zero.

\textbf{Factor analysis (H-E1):} Principal axis factoring on the 5×5 Spearman correlation matrix (factor_analyzer v0.4.x, Kaiser normalization). KMO adequacy threshold: > 0.60. Factor stability assessed via Tucker's congruence coefficient (φ ≥ 0.85).

\textbf{Capability independence (H-M1):} Survival fraction = |partial ρ| / |raw ρ|, threshold ≥ 0.50. Discriminant validity: |partial ρ(ECE, HumanEval|MMLU)| < 0.20.

\textbf{LOO adversarial prediction (H-M2):} Leave-one-out logistic regression (scikit-learn, C=1.0, max_iter=1000) with per-fold StandardScaler. Composite predictor (ECE + TruthfulQA\% + Brier) vs. MMLU-only baseline. ΔAUC with paired bootstrap CI (B=10,000). Gate thresholds: ΔAUC ≥ 0.10, CI lower bound > 0.

\textbf{Implementation:} scipy.stats.spearmanr for Spearman correlation; numpy seed 42 for reproducibility.

\subsubsection{3.5 Data Provenance}

The proof-of-concept uses \texttt{generate_synthetic_score_matrix()}, a parametric matrix generator with a pre-wired latent factor structure. This was detected by the pipeline's mock_data_check validator. The real-data pipeline (main.py + run_eval.py) is implemented using lm-evaluation-harness v0.4.x. Real-data execution requires approximately 2–4 GPU-hours per model.

